% CP-VI-0026
% target_chapter: VI/16 — Exact Sequences of Sheaves
% source_unit_id: IMP-DL-A99C7658F3-PSR000352-C80D6834
% source: Downloads/theory-of-algebraic-geometry-2.tex L352-L1174
% solution_provenance: SOURCE_IMPLICIT_SOLUTION

\subsection*{Problem 8. Quotients of Sheaves and Short Exact Sequences}

Let \(\mathcal{F}^{\prime}\) be a subsheaf of a sheaf \(\mathcal{F}\). We want to show that the natural map
\[
\mathcal{F}\longrightarrow \mathcal{F}/\mathcal{F}^{\prime}
\]
is surjective and has kernel \(\mathcal{F}^{\prime}\), so that there is an exact sequence
\[
0\longrightarrow \mathcal{F}^{\prime} \longrightarrow \mathcal{F} \longrightarrow \mathcal{F}/\mathcal{F}^{\prime} \longrightarrow 0.
\]

Then conversely, if
\[
0\longrightarrow \mathcal{F}^{\prime} \longrightarrow \mathcal{F} \longrightarrow \mathcal{F}^{\prime\prime} \longrightarrow 0
\]
is an exact sequence, we want to show that \(\mathcal{F}^{\prime}\) is isomorphic to a subsheaf of \(\mathcal{F}\), and \(\mathcal{F}^{\prime\prime}\) is isomorphic to the quotient of \(\mathcal{F}\) by this subsheaf.

\subsection{Part I: The quotient by a subsheaf gives a short exact sequence}

Let
\[
i:\mathcal{F}^{\prime}\hookrightarrow \mathcal{F}
\]
be the inclusion of a subsheaf.

The quotient presheaf is defined by
\[
U\longmapsto \mathcal{F}(U)/\mathcal{F}^{\prime}(U).
\]
However, this presheaf is not necessarily a sheaf. Therefore the quotient sheaf
\[
\mathcal{F}/\mathcal{F}^{\prime}
\]
is defined as the sheafification of this quotient presheaf.

There is a natural morphism
\[
q:\mathcal{F}\longrightarrow \mathcal{F}/\mathcal{F}^{\prime}.
\]

On stalks, quotients behave as expected:
\[
(\mathcal{F}/\mathcal{F}^{\prime})_x \cong \mathcal{F}_x/\mathcal{F}^{\prime}_x.
\]
Under this identification, the stalk map
\[
q_x:\mathcal{F}_x\longrightarrow (\mathcal{F}/\mathcal{F}^{\prime})_x
\]
is the ordinary quotient map
\[
\mathcal{F}_x\longrightarrow \mathcal{F}_x/\mathcal{F}^{\prime}_x.
\]

Therefore \(q_x\) is surjective for every \(x\in X\). Hence the morphism of sheaves
\[
q:\mathcal{F}\longrightarrow \mathcal{F}/\mathcal{F}^{\prime}
\]
is surjective.

Moreover, the kernel of \(q_x\) is
\[
\ker(q_x)=\mathcal{F}^{\prime}_x.
\]
Since kernels of morphisms of sheaves are determined stalkwise, we obtain
\[
\ker(q)=\mathcal{F}^{\prime}.
\]

Therefore the sequence
\[
0\longrightarrow \mathcal{F}^{\prime}
\xrightarrow{i}
\mathcal{F}
\xrightarrow{q}
\mathcal{F}/\mathcal{F}^{\prime}
\longrightarrow 0
\]
is exact.

Hence
\[
\boxed{
	0\longrightarrow \mathcal{F}^{\prime}
	\longrightarrow \mathcal{F}
	\longrightarrow \mathcal{F}/\mathcal{F}^{\prime}
	\longrightarrow 0
}
\]
is a short exact sequence of sheaves.

\subsection{Part II: Conversely, every short exact sequence is a quotient sequence}

Suppose that
\[
0\longrightarrow \mathcal{F}^{\prime}
\xrightarrow{i}
\mathcal{F}
\xrightarrow{p}
\mathcal{F}^{\prime\prime}
\longrightarrow 0
\]
is an exact sequence of sheaves.

Since the sequence is exact, the morphism
\[
i:\mathcal{F}^{\prime}\longrightarrow \mathcal{F}
\]
is injective. Hence \(\mathcal{F}^{\prime}\) is isomorphic to its image sheaf
\[
\operatorname{im}(i)\subseteq \mathcal{F}.
\]
Let
\[
\mathcal I=\operatorname{im}(i).
\]
Then \(\mathcal I\) is a subsheaf of \(\mathcal{F}\), and
\[
\mathcal{F}^{\prime}\cong \mathcal I.
\]

Because the sequence is exact, we have
\[
p\circ i=0.
\]
Therefore \(p\) vanishes on the subsheaf \(\mathcal I\). Hence \(p\) factors through the quotient sheaf
\[
\mathcal{F}/\mathcal I.
\]
Thus there exists an induced morphism
\[
\alpha:\mathcal{F}/\mathcal I\longrightarrow \mathcal{F}^{\prime\prime}
\]
such that the following diagram commutes:
\[
\begin{array}{ccccccccc}
	0 & \longrightarrow & \mathcal I & \longrightarrow & \mathcal{F} & \longrightarrow & \mathcal{F}/\mathcal I & \longrightarrow & 0 \\
	&                 & \downarrow  &               & \Vert &                 & \downarrow \alpha &                 &   \\
	0 & \longrightarrow & \mathcal{F}^{\prime} & \xrightarrow{i} & \mathcal{F} & \xrightarrow{p} & \mathcal{F}^{\prime\prime} & \longrightarrow & 0.
\end{array}
\]

We show that \(\alpha\) is an isomorphism.

Passing to stalks, we get
\[
\alpha_x:\mathcal{F}_x/\mathcal I_x\longrightarrow \mathcal{F}^{\prime\prime}_x.
\]
Since
\[
\mathcal I=\operatorname{im}(i),
\]
we have
\[
\mathcal I_x=\operatorname{im}(i_x).
\]
By exactness of the original sequence on stalks,
\[
\operatorname{im}(i_x)=\ker(p_x).
\]
Therefore
\[
\mathcal I_x=\ker(p_x).
\]

Also, since the sequence is exact, the morphism
\[
p_x:\mathcal{F}_x\longrightarrow \mathcal{F}^{\prime\prime}_x
\]
is surjective. Hence, by the first isomorphism theorem for abelian groups,
\[
\mathcal{F}_x/\ker(p_x)\cong \mathcal{F}^{\prime\prime}_x.
\]
But \(\ker(p_x)=\mathcal I_x\), so
\[
\mathcal{F}_x/\mathcal I_x\cong \mathcal{F}^{\prime\prime}_x.
\]
Thus
\[
\alpha_x:\mathcal{F}_x/\mathcal I_x\longrightarrow \mathcal{F}^{\prime\prime}_x
\]
is an isomorphism for every \(x\in X\).

Since a morphism of sheaves is an isomorphism if and only if it is an isomorphism on all stalks, it follows that
\[
\alpha:\mathcal{F}/\mathcal I\longrightarrow \mathcal{F}^{\prime\prime}
\]
is an isomorphism.

Therefore
\[
\mathcal{F}^{\prime\prime}\cong \mathcal{F}/\mathcal I.
\]

Since \(\mathcal I\cong \mathcal{F}^{\prime}\), we conclude that \(\mathcal{F}^{\prime}\) is isomorphic to a subsheaf of \(\mathcal{F}\), and \(\mathcal{F}^{\prime\prime}\) is isomorphic to the quotient of \(\mathcal{F}\) by that subsheaf.

Thus
\[
\boxed{
	\mathcal{F}^{\prime}\cong \mathcal I\subseteq \mathcal{F},
	\qquad
	\mathcal{F}^{\prime\prime}\cong \mathcal{F}/\mathcal I.
}
\]

\section{Theory}

Sheaf theory studies objects that are local in nature. A sheaf is not just one algebraic object, such as one abelian group. Instead, it assigns an abelian group to every open set of a topological space.

Let \(X\) be a topological space. A sheaf \(\mathcal{F}\) of abelian groups on \(X\) assigns to every open set \(U\subseteq X\) an abelian group
\[
\mathcal{F}(U),
\]
whose elements are called sections over \(U\).

If \(V\subseteq U\), there is a restriction homomorphism
\[
\rho^U_V:\mathcal{F}(U)\longrightarrow \mathcal{F}(V).
\]
For a section \(s\in \mathcal{F}(U)\), its restriction to \(V\) is usually denoted
\[
s|_V.
\]

The central idea is:

\[
\boxed{\text{Sheaves encode local-to-global information.}}
\]

Sections are local objects, but the sheaf axioms tell us when local pieces glue to a global object.

\subsection{The Sheaf Axioms}

\begin{definition}[Sheaf of abelian groups]
	A sheaf of abelian groups \(\mathcal{F}\) on a topological space \(X\) consists of:
	\begin{enumerate}[label=\textup{(\roman*)}]
		\item an abelian group \(\mathcal{F}(U)\) for each open set \(U\subseteq X\);
		\item restriction maps
		\[
		\rho^U_V:\mathcal{F}(U)\to \mathcal{F}(V)
		\]
		for each inclusion \(V\subseteq U\);
		\item the following compatibility conditions:
		\[
		\rho^U_U=\operatorname{id}_{\mathcal{F}(U)},
		\]
		and if \(W\subseteq V\subseteq U\), then
		\[
		\rho^U_W=\rho^V_W\circ \rho^U_V.
		\]
	\end{enumerate}
	Moreover, the following two sheaf axioms must hold.
	
	\textbf{Locality.} If \(\{U_i\}_{i\in I}\) is an open cover of \(U\), and if \(s,t\in \mathcal{F}(U)\) satisfy
	\[
	s|_{U_i}=t|_{U_i}
	\]
	for all \(i\in I\), then
	\[
	s=t.
	\]
	
	\textbf{Gluing.} If \(\{U_i\}_{i\in I}\) is an open cover of \(U\), and if sections
	\[
	s_i\in \mathcal{F}(U_i)
	\]
	satisfy
	\[
	s_i|_{U_i\cap U_j}=s_j|_{U_i\cap U_j}
	\]
	for all \(i,j\), then there exists a unique section
	\[
	s\in \mathcal{F}(U)
	\]
	such that
	\[
	s|_{U_i}=s_i
	\]
	for all \(i\).
\end{definition}

\subsection{Stalks and Germs}

The most important local object associated with a sheaf is its stalk.

\begin{definition}[Stalk]
	Let \(\mathcal{F}\) be a sheaf on \(X\), and let \(x\in X\). The stalk of \(\mathcal{F}\) at \(x\) is
	\[
	\mathcal{F}_x=\operatorname*{colim}_{x\in U}\mathcal{F}(U),
	\]
	where the colimit is taken over all open neighbourhoods \(U\) of \(x\).
\end{definition}

An element of \(\mathcal{F}_x\) is called a germ.

If \(s\in \mathcal{F}(U)\) and \(x\in U\), then the image of \(s\) in \(\mathcal{F}_x\) is denoted
\[
s_x.
\]

Two sections \(s\in \mathcal{F}(U)\) and \(t\in \mathcal{F}(V)\) define the same germ at \(x\) if there exists an open neighbourhood
\[
W\subseteq U\cap V
\]
of \(x\) such that
\[
s|_W=t|_W.
\]

Thus the stalk \(\mathcal{F}_x\) remembers only what happens arbitrarily close to \(x\).

\begin{remark}
	The stalk is the local algebraic object attached to the point \(x\). It ignores global behaviour and remembers only local behaviour near \(x\).
\end{remark}

\subsection{Morphisms of Sheaves}

\begin{definition}[Morphism of sheaves]
	Let \(\mathcal{F}\) and \(\mathcal{G}\) be sheaves of abelian groups on \(X\). A morphism of sheaves
	\[
	\varphi:\mathcal{F}\longrightarrow \mathcal{G}
	\]
	is a family of group homomorphisms
	\[
	\varphi_U:\mathcal{F}(U)\longrightarrow \mathcal{G}(U)
	\]
	for every open set \(U\subseteq X\), compatible with restrictions.
	
	That is, whenever \(V\subseteq U\), the diagram
	\[
	\begin{array}{ccc}
		\mathcal{F}(U) & \xrightarrow{\varphi_U} & \mathcal{G}(U) \\
		\downarrow & & \downarrow \\
		\mathcal{F}(V) & \xrightarrow{\varphi_V} & \mathcal{G}(V)
	\end{array}
	\]
	commutes.
\end{definition}

In other words, for every \(s\in \mathcal{F}(U)\),
\[
\varphi_U(s)|_V=\varphi_V(s|_V).
\]

Usually we simply write
\[
\varphi(s)|_V=\varphi(s|_V).
\]

A morphism of sheaves induces a morphism on stalks. For each \(x\in X\), we get
\[
\varphi_x:\mathcal{F}_x\longrightarrow \mathcal{G}_x
\]
defined by
\[
\varphi_x(s_x)=(\varphi(s))_x.
\]

This is well-defined because \(\varphi\) is compatible with restrictions.

\subsection{Kernels, Images, and Cokernels}

\subsubsection{Kernel sheaf}

Let
\[
\varphi:\mathcal{F}\to \mathcal{G}
\]
be a morphism of sheaves.

The kernel is defined sectionwise:
\[
(\ker\varphi)(U)=\ker\bigl(\mathcal{F}(U)\to \mathcal{G}(U)\bigr).
\]

This is already a sheaf.

\begin{proposition}
	For every point \(x\in X\),
	\[
	(\ker\varphi)_x=\ker(\varphi_x).
	\]
\end{proposition}

\subsubsection{Image sheaf}

The image is more subtle.

One may first define the image presheaf by
\[
U\longmapsto \operatorname{im}\bigl(\mathcal{F}(U)\to \mathcal{G}(U)\bigr).
\]

However, this presheaf need not be a sheaf. Therefore the image sheaf is defined as the sheafification of this presheaf:
\[
\operatorname{im}(\varphi)=
\left(
U\longmapsto \operatorname{im}(\varphi_U)
\right)^{\#}.
\]

On stalks, nevertheless, the image behaves exactly as expected:
\[
(\operatorname{im}\varphi)_x=\operatorname{im}(\varphi_x).
\]

\subsubsection{Cokernel sheaf}

Similarly, the cokernel presheaf is
\[
U\longmapsto \operatorname{coker}\bigl(\mathcal{F}(U)\to \mathcal{G}(U)\bigr).
\]

The cokernel sheaf is its sheafification:
\[
\operatorname{coker}(\varphi)=
\left(
U\longmapsto \operatorname{coker}(\varphi_U)
\right)^{\#}.
\]

On stalks,
\[
(\operatorname{coker}\varphi)_x=\operatorname{coker}(\varphi_x).
\]

\subsection{Exactness of Sheaves}

\begin{definition}[Exact sequence of sheaves]
	A sequence of sheaves
	\[
	A\longrightarrow B\longrightarrow C
	\]
	is exact at \(B\) if
	\[
	\operatorname{im}(A\to B)=\ker(B\to C).
	\]
\end{definition}

The crucial theorem is that exactness of sheaves can be checked on stalks.

\begin{theorem}[Exactness is stalkwise]
	A sequence of sheaves
	\[
	A\longrightarrow B\longrightarrow C
	\]
	is exact if and only if for every point \(x\in X\), the induced sequence of stalks
	\[
	A_x\longrightarrow B_x\longrightarrow C_x
	\]
	is exact as a sequence of abelian groups.
\end{theorem}

\begin{remark}
	This is one of the main reasons stalks are so important. Sheaf exactness is a local condition.
\end{remark}

Thus, to prove exactness of sheaves, it is usually enough to prove ordinary exactness of abelian groups after passing to stalks.

\subsection{Injective and Surjective Morphisms}

\subsubsection{Injective morphisms}

A morphism
\[
\varphi:\mathcal{F}\to \mathcal{G}
\]
is injective if
\[
\ker\varphi=0.
\]

Equivalently, for every \(x\in X\),
\[
\varphi_x:\mathcal{F}_x\to \mathcal{G}_x
\]
is injective.

If \(\varphi\) is injective and if \(s,t\in \mathcal{F}(U)\) satisfy
\[
\varphi(s)=\varphi(t),
\]
then
\[
s=t.
\]

Thus injectivity behaves very similarly to ordinary injectivity.

\subsubsection{Surjective morphisms}

A morphism
\[
\varphi:\mathcal{F}\to \mathcal{G}
\]
is surjective if for every \(x\in X\),
\[
\varphi_x:\mathcal{F}_x\to \mathcal{G}_x
\]
is surjective.

This means that every germ of a section of \(\mathcal{G}\) locally comes from a section of \(\mathcal{F}\).

In other words, if \(t\in \mathcal{G}(U)\) and \(x\in U\), then there exists a neighbourhood
\[
V\subseteq U
\]
of \(x\), and a section
\[
s\in \mathcal{F}(V),
\]
such that
\[
\varphi(s)=t|_V.
\]

So surjectivity of sheaves means local liftability.

It does not necessarily mean that for every open set \(U\), the map
\[
\mathcal{F}(U)\to \mathcal{G}(U)
\]
is surjective.

This distinction is fundamental.

\[
\boxed{
	\text{Surjective as a sheaf means locally surjective, not necessarily globally surjective.}
}
\]

\subsection{Isomorphisms of Sheaves}

\begin{theorem}
	A morphism of sheaves
	\[
	\varphi:\mathcal{F}\to \mathcal{G}
	\]
	is an isomorphism if and only if it is injective and surjective.
\end{theorem}
